%% kapa: exam=1 round=28 year=2017 period=1 form=A subject=civil_law subject_ko=민법 seq=17 number=17
%% source: Q-Net 감정평가사 기출문제 (cst00309) artlSeq=5199603 file=extracted/1008/5199603/2017년 제28회 감정평가사 1차시험 문제지/1교시 A형.pdf
%% conversion: kordoc 4.24.51 -> tools/gen1.py (auto); review=proofread-sonnet; flags=none
\begin{question}{17}
소멸시효 완성 후에 한 시효이익의 포기에 관한 설명으로 옳지 않은 것은? (다툼이 있으면 판례에 따름)
\begin{choices}
  \item 시효이익을 포기하면 그 때부터 시효가 새로 진행한다.
  \item 시효완성의 이익을 받을 당사자 또는 그 대리인은 시효이익 포기의 의사표시를 할 수 있다.
  \item 주채무자가 시효이익을 포기하더라도 보증인에게는 그 효력이 없다.
  \item 시효이익을 이미 포기한 사람과의 법률관계를 통해 시효이익을 원용할 이해관계를 형성한 사람은 소멸시효를 주장할 수 있다.
  \item 채권의 시효완성 후에 채무자가 그 기한의 유예를 요청한 때에는 시효이익을 포기한 것으로 보아야 한다.
\end{choices}
\end{question}
\answer{4}
